% Source: https://artofproblemsolving.com/community/c3h1526212_last_post_continued
% Stable identity: aops-topic-1526212-q11
A positive integer has an odd number of divisors exactly when it is a perfect square (divisors pair up as $d$ and $\dfrac {n}{d}$, except $\sqrt n$). On a die the squares are $1$ and $4$, so she steps right with probability $\dfrac{2}{6}=\dfrac{1}{3}$ and left with probability $\dfrac{2}{3}$.

She is back at the start after four steps exactly when two steps go right and two go left:
\[
\dbinom42\left(\dfrac{1}{3}\right)^2\left(\dfrac{2}{3}\right)^2=6\cdot\dfrac{1}{9}\cdot\dfrac{4}{9}=\boxed{\dfrac{8}{27}}.
\]
